\section{The chromatic expansion}
\label{main:completion}

We return to the original graph and its graph-nondescent permutations.

\begin{proof}[Proof of Theorem~\ref{main:theorem}]
Choose the triangular embedding of Lemma~\ref{tri:embedding}, and
fix the rational generic choices and finite local orders in
Section~\ref{route:section} by the enumeration specified there.
Theorem~\ref{route:witness} gives a terminating map
$\Lambda_G:\mathcal A_G\to\{\lambda:\lambda\vdash n\}$ with
\[
 \sum_{\substack{w\in\mathcal A_G\\\Lambda_G(w)=\lambda}}
       q^{\inv_G(w)}
   =[e_\lambda]X_G(X;q)=c_\lambda(q).
\]
For $\sigma\in D_G^0$ define
\[
 \theta_G(\sigma)=\Lambda_G(\operatorname{reverse}\sigma).
\]
By Lemma~\ref{graph:digits}, reversal gives the required input to
$\Lambda_G$ and complements graph inversions. Lemma~\ref{graph:reciprocity}
then gives, for every $\lambda$,
\[
 \sum_{\substack{\sigma\in D_G^0\\\theta_G(\sigma)=\lambda}}
       q^{\ginv_G(\sigma)}
   =q^M c_\lambda(q^{-1})=c_\lambda(q).
\]
Together with $X_G=\chi_G$, this proves
Equation~\eqref{main:witness}. Every coefficient is the generating
polynomial of a finite set, so it belongs to $\N[q]$. All choices and
all recursions are finite procedures for each input. For $n=0$, take
the empty graph, the empty permutation and the empty partition.
\end{proof}

\begin{remark}[Scope of the algorithm]
The map uses the graph, the input permutation, finite wall recursions
and local coefficient indexings. The elementary coefficient list is
not supplied to it. The result proves termination, but gives no
polynomial-time bound. The construction is a mathematical algorithm
based on finite rational arithmetic and ordered finite sets.
\end{remark}
